% SPDX-License-Identifier: Apache-2.0
\section{What Was Compiled}
\label{sec:e-probes}

Every claim in this paper about what Rust does is a probe that was built.
The probes are one file per question in \code{experiments/rust\_embedding},
and those expected to fail are tagged \code{manual} and stay out of the
default build; the \code{BUILD.bazel} beside them lists which, each with
the error it gives. A failing probe is kept rather than deleted,
because the failure is the answer.

\begin{table*}[t]
\caption{The probes and what they answered. Rust 1.90.0 stable unless
stated. Two results changed the design rather than confirming it.}
\label{tab:probes}
\centering
\footnotesize
\begin{tabular}{@{}lp{60mm}p{80mm}@{}}
\toprule
Probe & Question & Answer \\
\midrule
\code{probe\_widths}
  & Does \code{U<\{A + B\}>} compile?
  & \textbf{No.} \code{error: generic parameters may not be used in const operations} \\
\code{probe\_widths\_nightly}
  & Does it compile on nightly?
  & \textbf{Yes}, with \code{generic\_const\_exprs}, and \code{mul(a, b)} resolves to \code{U<64>} at the call site \\
\code{probe\_design}
  & Is \code{async fn} in a trait usable?
  & \textbf{Yes}, warning that auto trait bounds such as \code{Send} cannot be stated \\
\code{probe\_design}
  & Inputs and outputs as type parameters?
  & \textbf{Yes.} One struct implements \code{Unit} twice with different shapes \\
\code{probe\_design}
  & Interface as struct, modport as trait?
  & \textbf{Yes}, including a unit generic over the role rather than the bus \\
\code{probe\_comp}
  & Do \code{Bus}, \code{Port}, \code{Signal}, \code{Chan} hold together?
  & \textbf{Yes}, and construction compiles \\
\code{probe\_comp}
  & Is a \code{Tag} type parameter workable?
  & \textbf{Yes}, with the policy as associated constants readable in the body \\
\code{probe\_attr\_block}
  & Is \code{\#[tag(..)]} on a block stable?
  & \textbf{No.} \code{error[E0658]: attributes on expressions are experimental} \\
\code{probe\_processes}
  & Can \code{run} join two processes?
  & \textbf{Yes}, with processes taking \code{\&self} and registers as cells \\
\code{probe\_processes\_mut}
  & Can they take \code{\&mut self}?
  & \textbf{No.} \code{error[E0499]: cannot borrow *self as mutable more than once at a time} \\
\code{probe\_config}
  & Is a config a struct with associated types and constants?
  & \textbf{Yes}, and its constants work in an array length, so they are genuinely compile-time \\
\code{probe\_item\_order}
  & May an \code{impl} precede the \code{struct} it is for?
  & \textbf{Yes.} Items in a module are not order dependent, including generically and for types declared later \\
\code{probe\_derive}
  & Does a derive give the same reading and stay idiomatic?
  & \textbf{Yes}, in about thirty lines needing neither \code{syn} nor \code{quote} \\
\code{probe\_funcs}
  & Can operator latency replace a hand-placed stage boundary?
  & \textbf{Yes.} An \code{async} operator in \code{crate::pipeline} is awaited, and the depth follows from the operators used \\
\code{probe\_if}
  & Do \code{when!} and \code{with!} work, the fields as entries, the condition first or the struct first?
  & \textbf{Yes}, as procedural macros; \code{macro\_rules!} cannot follow an \code{expr} with \code{?} or \code{<=} \\
\code{probe\_when\_bad}
  & An entry with no colon?
  & \textbf{Refused:} \code{expected `field: value`, `c ? field: value` or `c ? \{ .. \}`} \\
\code{probe\_case}
  & Does \code{case!} take Rust patterns, with guards and alternatives?
  & \textbf{Yes}, tested with \code{matches!}; the first arm that matches wins \\
\code{probe\_case\_bad}
  & A drive with no arrow in an arm?
  & \textbf{Refused:} \code{expected `register <= value`}, the macro's own message \\
\code{probe\_case\_binding}
  & May an arm's pattern bind, as \code{Op::Load(v)}?
  & \textbf{No.} \code{error[E0425]: cannot find value `v`}; \code{matches!} scopes its bindings \\
\code{probe\_slice}
  & May a slice bound be a const parameter, or a run-time offset?
  & \textbf{Yes} to both, with the width static \\
\code{probe\_slice\_dyn}
  & May the slice \emph{width} be a run-time value?
  & \textbf{No.} \code{error[E0435]: attempt to use a non-constant value in a constant} \\
\code{probe\_hierarchy}
  & How is a unit with submodules instantiated and wired?
  & Submodules and their buses are fields; ports are passed to \code{run}. Disjoint fields, so the borrow rule of probe 5b does not apply \\
\code{probe\_main}
  & Can \code{fn main()} instantiate the top through a config?
  & \textbf{Yes}, and it runs: \code{bazel run ... -- asic} prints \code{elaborated config 'asic' at 400 MHz} \\
\code{probe\_ports}
  & Can direction be enforced rather than documented?
  & \textbf{Yes}, by which methods an end exposes \\
\code{probe\_ports\_bad}
  & Can a unit drive an input?
  & \textbf{No.} \code{error[E0599]: no method named `set` found for struct `In`} \\
\code{probe\_values}
  & Do \code{Bit}, \code{Logic}, \code{U<N>}, \code{I<N>} work on stable?
  & \textbf{Yes}, including the IEEE 1164 resolution table and \code{U} as the default of an undriven wire \\
\code{probe\_signal\_split}
  & Can a signal yield its two ends separately, as \code{mpsc} does?
  & \textbf{Yes}, with \code{In} clonable for fanout and \code{Out} unique \\
\code{probe\_two\_drivers}
  & Can a wire have two drivers?
  & \textbf{No.} \code{error[E0382]: use of moved value} \\
\code{probe\_nested\_config}
  & Can a submodule have a configuration of its own?
  & \textbf{Yes.} A parent's config names its children's, so two children may share a knob name and a child's config is reusable \\
\code{probe\_interface\_macro}
  & Can \code{macro\_rules!} generate the role views?
  & \textbf{Yes}, for exactly two roles, spelling \code{in} as \code{inp}, and a role that omits a member compiles silently \\
\code{probe\_interface\_proc}
  & Can a procedural macro do it properly?
  & \textbf{Yes.} Any number of roles, \code{in} as written, under two hundred lines, no \code{syn} \\
\code{probe\_interface\_incomplete}
  & A role that omits a member?
  & \textbf{Refused:} \code{role `Target` does not say which end of `dat` it takes} \\
\code{probe\_interface\_twodrivers}
  & Two roles driving one member?
  & \textbf{Refused:} \code{member `adr` is driven by more than one role: A, B} \\
\code{probe\_domain}
  & Can the clock domain be a type parameter with a named default?
  & \textbf{Yes.} A single-clock design never mentions one; a two-clock design names the second \\
\code{probe\_domain\_cross}
  & A wire across clocks with nothing between?
  & \textbf{Refused:} \code{expected In<u32, Clk400>, found In<u32, Clk100>}, and the default is not a wildcard \\
\code{probe\_infer}
  & May a const width be left to inference, as \code{\_}?
  & \textbf{Yes}, since Rust 1.89, where the context fixes it: an annotated destination, a typed use, the other operand, the return type \\
\code{probe\_infer\_bad}
  & A width nothing fixes, a slice compared with a number?
  & \textbf{No.} \code{error[E0284]: type annotations needed}; and the lowering, which sees no types, asks for every width it needs \\
\code{probe\_infer\_lower}
  & A width Rust would infer, under \code{\#[lower]}?
  & \textbf{Refused:} \code{`slice::<..>` needs its width written: the lowering sees no types}; \code{concat}'s low operand is the one width it does not need \\
\code{probe\_macro\_in\_macro}
  & Does \code{\#[lower]} written by a macro's output expand, and does a port's field, \code{p.tag}, lower?
  & \textbf{Yes} to both: the unit lowers, and the field is the slice of the port's data the value's layout gives it \\
\bottomrule
\end{tabular}
\end{table*}

\subsection{The four that changed the design}

\textbf{Arbitrary-width arithmetic needs nightly.} A hardware library wants
to write that multiplying an $A$-bit value by a $B$-bit value gives an
$A+B$-bit one. Stating that in a signature needs arithmetic over const
generic parameters, and stable Rust refuses it. The choices are nightly
with \code{generic\_const\_exprs}, or one function per width pair, which
compiles on stable and does not scale. This is the largest single
constraint on the proposal, and it was a guess until it was built.

\textbf{Parallel processes cannot take a mutable borrow.}
Section~\ref{sec:e-units} gives the collision and the repair. The point
here is that ten lines of Rust found a constraint that reading the language
reference had not suggested, and that the repair, a register being a cell
several processes reach, is a better model of a register than the one that
failed.

\textbf{A \code{macro\_rules!} conditional cannot reach inside a block.}
Writing \code{with!} showed where the macro route stops, and that is the
strongest argument for the recommendation in
Section~\ref{sec:e-elaboration}. Section~\ref{sec:e-walls} states it.

\textbf{One driver per wire comes free.} It was not designed in.
\code{signal} returns one \code{Out} and \code{Out} is not \code{Clone},
so a second driver is a value that cannot be made, and the oldest error
message in hardware design arrives as \code{error[E0382]: use of moved
value}. Section~\ref{sec:e-wiring} states the cardinality this rests on.

\subsection{On method}

The first draft of this work argued from the language rules and was wrong
twice in eleven claims. Both errors were in the direction of optimism, and
both were cheap to find. A design document for an embedded language should
include its probes, and the probes should be part of the build, so that a
Rust release that changes an answer breaks something visible rather than
quietly invalidating a paragraph.
